\documentclass{article}
\usepackage{iclr2027_conference,times}
\usepackage[utf8]{inputenc}
\usepackage[T1]{fontenc}
\usepackage[italian]{babel}
\usepackage{amsmath,amssymb,amsfonts}
\usepackage{booktabs}
\usepackage{graphicx}
\usepackage{xcolor}
\usepackage{hyperref}
\usepackage{url}
\usepackage{microtype}
\usepackage{algorithm}
\usepackage{algpseudocode}

\title{Propagazione di survey termodinamica\\
\large Deformazione di semianello e controllo della complessit\`a nel $K$-SAT casuale}

\author{Carlo Nicolini\\
Raffaele Marino\\
\texttt{c.nicolini@ipazia.com}}

\newcommand{\E}{\mathbb{E}}
\newcommand{\Sig}{\Sigma}
\newcommand{\Tc}{T_{\mathrm{cav}}}
\newcommand{\Ta}{T_{\mathrm{act}}}

\begin{document}
\maketitle

\begin{abstract}
Nel $K$-SAT casuale, vicino alla soglia di soddisfacibilit\`a, lo spazio delle soluzioni si frammenta in un numero esponenziale di cluster.
La survey propagation (SP) stima la complessit\`a $\Sig$, logaritmo del numero di cluster, e guida la decimazione.
Il backtracking survey propagation (BSP) corregge scelte precedenti per mantenere $\Sig$ positiva.
In questo lavoro proponiamo una lettura termodinamica di SP: l'aggregazione ``dura'' delle configurazioni di warning viene sostituita dall'operazione di free energy dei semianelli termodinamici di Marcolli e Thorngren, con temperatura cavitarie $\Tc$.
Nel limite $\Tc\to 0$ si recupera SP; a $\Tc>0$ la costruzione coincide con la famiglia SP a energia finita, reinterpretata come deformazione di semianello.
Sullo stesso asse costruiamo una policy di decisione basata su risposte locali della free entropy, di cui un membro a vincolo duro \`e lo stimatore $I(k)$ della frazione di cluster compatibili con un'assegnazione gi\`a fissata.
Su istanze $3$-SAT e $4$-SAT, la policy $I(k)$ predice la variazione reale di $\Sig$ al rilascio con correlazione $r\simeq 0{,}997$ e migliora la dinamica della complessit\`a rispetto ai criteri di certainty e polarizzazione.
Restano aperti lo spostamento della soglia algoritmica $\alpha_a$ e la costruzione operativa del bias verso cluster non congelati.
\end{abstract}

\section{Introduzione}
\label{sec:intro}

Il problema di soddisfacibilit\`a booleana casuale ($K$-SAT) \`e un banco di prova per la meccanica statistica dei sistemi disordinati e per gli algoritmi di ottimizzazione discreta~\cite{mezard2009information}.
Una formula con $N$ variabili e $M=\alpha N$ clausole di lunghezza $K$ definisce un paesaggio di energia $E(\sigma)$ uguale al numero di clausole violate.
A temperatura zero le configurazioni di energia nulla sono le soluzioni.
Al crescere di $\alpha$ lo spazio delle soluzioni cambia di natura: oltre una densit\`a dinamica $\alpha_d$ si spezza in un numero esponenziale di cluster, e oltre una densit\`a di condensazione $\alpha_c$ un numero sotto-esponenziale di cluster porta quasi tutta la misura~\cite{krzakala2007clusters,montanari2008clusters}.
La soglia di soddisfacibilit\`a $\alpha_s$ segna la scomparsa delle soluzioni~\cite{mezard2002analytic,achlioptas2005rigidity}.

La teoria della rottura di simmetria di replica a un passo (one-step replica symmetry breaking, 1RSB) descrive questa geometria~\cite{parisi1980sequence,mezard2001bethe}.
I cluster sono stati puri della misura di Gibbs; due soluzioni dello stesso cluster hanno overlap alto, due soluzioni di cluster diversi overlap basso.
Il parametro di Parisi $m$ pesa i cluster secondo la loro entropia interna.
Survey propagation corrisponde al conteggio uniforme dei cluster, cio\`e al limite $m\to 0$ della 1RSB a temperatura fisica zero~\cite{mezard2002analytic,braunstein2005survey}.
Le equazioni SP possono anche essere rilette come belief propagation su uno spazio esteso in cui ogni variabile ammette lo stato ``*''~\cite{braunstein2004surveybp}.

Gli algoritmi basati su SP fissano progressivamente le variabili pi\`u polarizzate (decimazione) e aggiornano i messaggi sulla formula residua.
Il BSP di Marino, Parisi e Ricci-Tersenghi~\cite{marino2016bsp} aggiunge mosse di backtracking: variabili gi\`a fissate possono essere liberate, con frequenza regolata da un rapporto $r$ fra backtracking e decimazione.
L'obiettivo dichiarato \`e mantenere la complessit\`a $\Sig$ positiva il pi\`u a lungo possibile, perch\'e una $\Sig$ che crolla segnala la perdita di troppi cluster e spesso precede il fallimento della ricerca.
Gli autori stessi osservano che una misura concentrata sui cluster non congelati (unfrozen) potrebbe ridurre il backtracking necessario, ma non forniscono un criterio locale operativo per seguirla.

Due lacune restano aperte.
La prima riguarda la metrica di decisione: BSP ordina decimazione e backtracking con bias locali (polarizzazione, certainty) calcolati al momento della scelta.
Per le variabili fissate, i bias non vengono aggiornati con i messaggi correnti, e il confronto fra scelte di et\`a diverse mescola quantit\`a definite su formule residue diverse.
La seconda riguarda la struttura dell'aggregazione cavitarie: SP esclude in modo duro ogni configurazione di warning con conflitti.
Esiste gi\`a una letteratura di SP a energia finita, ma manca una lettura che colleghi in modo esplicito quella deformazione ai semianelli termodinamici~\cite{marcolli2014thermodynamic}, costruiti sulla teoria di Connes e Consani in caratteristica uno~\cite{connes2010characteristic,connes2014witt}, e a una policy di controllo del processo di ricerca~\cite{kappen2005path}.

In questo lavoro sviluppiamo entrambe le linee.
Nella Sezione~\ref{sec:theory} richiamiamo 1RSB, SP e BSP.
Nella Sezione~\ref{sec:thermo} definiamo la survey propagation termodinamica (ThermoSP): la somma sulle configurazioni di warning viene sostituita dall'operazione di free energy di Shannon a temperatura $\Tc$, e dimostriamo il recupero di SP per $\Tc\to 0$.
Nella Sezione~\ref{sec:policy} introduciamo politiche basate su risposte della free entropy, incluso lo stimatore $I(k)$ delle survey virtuali.
Nella Sezione~\ref{sec:methods} descriviamo l'implementazione e i protocolli numerici.
Nella Sezione~\ref{sec:results} riportiamo i risultati preliminari.
Nella Sezione~\ref{sec:discussion} discutiamo limiti e prospettive, e nella Sezione~\ref{sec:conclusions} concludiamo.

\section{Quadro teorico}
\label{sec:theory}

\subsection{Complessit\`a e misura 1RSB}

Sia $\mathcal{C}$ l'insieme dei cluster di soluzioni di una formula fissata.
Indichiamo con $s_C=N^{-1}\log|C|$ l'entropia interna del cluster $C$ e con $\Sig(f)$ la densit\`a di complessit\`a, cio\`e $N^{-1}$ volte il logaritmo del numero di cluster con densit\`a di energia libera $f$.
La somma pesata
\begin{equation}
  \Xi(m)=\sum_{C\in\mathcal{C}}|C|^m
  \simeq\int\mathrm{d}f\,\exp\bigl(N[\Sig(f)-\beta m f]\bigr)
\end{equation}
seleziona, al variare di $m$, diverse sezioni della curva $\Sig$.
A $m\to 0$ i cluster sono contati uno a uno; a $m=1$ si recupera l'equilibrio, in cui i cluster pesano secondo il numero di soluzioni che contengono.
Nel codice e negli esperimenti di questo lavoro la quantit\`a chiamata complessit\`a \`e il logaritmo estensivo del numero di cluster; la densit\`a \`e quindi $\Sig/N$.

\subsection{Equazioni di survey propagation}

Sul factor graph della formula, i messaggi di SP sono survey $\eta_{a\to i}\in[0,1]$: probabilit\`a che la clausola $a$ invii un warning alla variabile $i$ in un cluster tipico.
Per un arco variabile--clausola $j\to a$, i warning entranti dagli altri vicini di $j$ producono un campo cavitarie $h$.
Le configurazioni con conflitti (warning opposti) sono escluse.
Le quantit\`a $\Pi^u_{j\to a}$, $\Pi^s_{j\to a}$ e $\Pi^0_{j\to a}$ raccolgono le probabilit\`a delle configurazioni che spingono $j$ a non soddisfare $a$, a soddisfarla, o a lasciarla libera.
L'aggiornamento di clausola \`e
\begin{equation}
  \eta_{a\to i}
  =\prod_{j\in\partial a\setminus i}
  \frac{\Pi^u_{j\to a}}{\Pi^u_{j\to a}+\Pi^s_{j\to a}+\Pi^0_{j\to a}}.
  \label{eq:sp-eta}
\end{equation}
Dai messaggi convergenti si ottengono i bias di variabile $w_i^+$, $w_i^-$ e la complessit\`a $\Sig$ come somma di contributi di clausole e di variabili~\cite{braunstein2005survey,marino2016bsp}.

\subsection{Backtracking survey propagation}

BSP alterna aggiornamenti SP, mosse di decimazione e mosse di backtracking.
Nella decimazione si fissa la variabile con bias pi\`u alto nella direzione suggerita.
Nel backtracking si libera una variabile gi\`a fissata secondo un criterio locale (nel BSP originale: polarizzazione o certainty calcolate al momento del fissaggio).
Il parametro $r$ \`e il rapporto fra numero di backtracking e numero di decimazioni; per $r=0$ si recupera la decimazione pura guidata da SP.
Il tempo di calcolo scala come $N/(1-r)$~\cite{marino2016bsp}.

\section{Survey propagation termodinamica}
\label{sec:thermo}

\subsection{Semianelli termodinamici}

Nel semianello tropicale min-plus l'addizione \`e $x\oplus y=\min(x,y)$: sopravvive soltanto l'alternativa migliore.
Marcolli e Thorngren~\cite{marcolli2014thermodynamic} deformano questa operazione con un'entropia $S$ e una temperatura $T$:
\begin{equation}
  x\oplus_{S,T}y
  =\min_{p\in[0,1]}
  \bigl[p\,x+(1-p)\,y-T\,S(p)\bigr].
\end{equation}
Per l'entropia di Shannon si ottiene
\begin{equation}
  x\oplus_{\mathrm{Sh},T}y
  =-T\log\bigl(e^{-x/T}+e^{-y/T}\bigr),
\end{equation}
cio\`e una free energy a due stati.
Per $T\to 0$ si recupera $\min(x,y)$.
La stessa costruzione ammette una forma $n$-aria sul simplesso delle probabilit\`a, interpretabile come mescolamento termodinamico di $n$ alternative~\cite{marcolli2014thermodynamic}.
Questa \`e la struttura algebrica che usiamo per aggregare le configurazioni di warning.

\subsection{Deformazione delle somme cavitarie}

Fissiamo un arco $j\to a$ e una configurazione $\omega$ dei warning entranti in $j$ dagli altri vicini.
Siano $p(\omega)$ e $q(\omega)$ i numeri di warning che spingono $j$ a soddisfare $a$ e a non soddisfarla.
L'energia di conflitto \`e
\begin{equation}
  \Delta E(\omega)=\min\bigl(p(\omega),q(\omega)\bigr).
\end{equation}
Se $P_0(\omega)$ \`e la misura a priori indotta dalle survey entranti, definiamo la free energy del settore di campo $h$ come
\begin{equation}
  F^{(T)}_{j\to a}(h)
  =-T\log\sum_{\omega:\,h(\omega)=h}
  P_0(\omega)\,e^{-\Delta E(\omega)/T}.
  \label{eq:F-thermo}
\end{equation}
In forma di semianello,
\begin{equation}
  F^{(T)}_{j\to a}(h)
  =\bigoplus_{\omega:\,h(\omega)=h}^{\mathrm{Sh},T}
  \bigl[\Delta E(\omega)-T\log P_0(\omega)\bigr].
\end{equation}
La forma variazionale equivalente \`e
\begin{equation}
  F^{(T)}_{j\to a}(h)
  =\min_{\rho}
  \Bigl[
  \E_{\rho}[\Delta E]+T\,D_{\mathrm{KL}}(\rho\Vert P_0)
  \Bigr],
\end{equation}
dove il minimo corre sulle misure $\rho$ supportate sulle configurazioni con campo $h$.
La soluzione ``dura'' minimizza soltanto l'energia; la soluzione termodinamica bilancia energia e distanza informazionale dalla misura cavitaria di riferimento.

Le probabilit\`a termiche dei settori $u$, $s$ e $0$ sono
\begin{align}
  \widetilde{\Pi}^{u}_{j\to a}(T)
  &\propto\sum_{q>p}C_{pq}\,e^{-\min(p,q)/T},\\
  \widetilde{\Pi}^{s}_{j\to a}(T)
  &\propto\sum_{p>q}C_{pq}\,e^{-\min(p,q)/T},\\
  \widetilde{\Pi}^{0}_{j\to a}(T)
  &\propto\sum_{p=q}C_{pq}\,e^{-p/T},
\end{align}
dove $C_{pq}$ \`e la probabilit\`a a priori di $p$ warning soddisfacenti e $q$ insoddisfacenti, ottenuta dal prodotto polinomiale sulle survey entranti.
L'equazione di clausola diventa
\begin{equation}
  \eta^{(T)}_{a\to i}
  =\prod_{j\in\partial a\setminus i}
  \frac{\widetilde{\Pi}^{u}_{j\to a}(T)}
  {\widetilde{\Pi}^{u}_{j\to a}(T)
  +\widetilde{\Pi}^{s}_{j\to a}(T)
  +\widetilde{\Pi}^{0}_{j\to a}(T)}.
  \label{eq:thermo-eta}
\end{equation}

\subsection{Limiti e identificazioni}

Per $T\to 0$,
\begin{equation}
  e^{-\min(p,q)/T}
  \to
  \begin{cases}
  1 & \text{se }\min(p,q)=0,\\
  0 & \text{se }\min(p,q)>0.
  \end{cases}
\end{equation}
Le somme $\widetilde{\Pi}$ riducono alle $\Pi$ di SP, e quindi
\begin{equation}
  \lim_{T\to 0}\mathrm{ThermoSP}(T)=\mathrm{SP}.
\end{equation}
A $T>0$ il peso $e^{-\Delta E/T}$ \`e il fattore di Boltzmann delle configurazioni con conflitti.
Identificando $T=1/y$ si recupera la survey propagation a energia finita gi\`a presente in letteratura~\cite{braunstein2004surveybp,wemmenhove2007finite}.
La novit\`a che rivendichiamo non \`e quindi la sola deformazione dei messaggi, ma la sua lettura come operazione di semianello termodinamico e, soprattutto, l'uso delle free energy risultanti come segnali di controllo del BSP (Sezione~\ref{sec:policy}).

Due temperature devono restare distinte.
La temperatura cavitarie $\Tc$ deforma l'aggregazione dei warning dentro le equazioni di messaggio.
La temperatura di azione $\Ta$ deforma la policy di scelta fra azioni (decimare, backtrackare, scegliere il segno).
Sono parametri indipendenti: $\Tc$ ha un interpretazione fisica nel sistema esteso di~\cite{braunstein2004surveybp}; $\Ta$ \`e un parametro di esplorazione del solver.

\section{Politiche basate sulla free entropy}
\label{sec:policy}

\subsection{Dal bias al costo di free energy}

Nel BSP standard i bias $w_i^\pm$ sono marginali locali.
Noi consideriamo invece risposte della free entropy a un vincolo locale.
Siano $\Phi_i^+(T)$ e $\Phi_i^-(T)$ le free energy cavitarie associate ai due valori di $x_i$.
Il bias termodinamico \`e
\begin{equation}
  B_i^{\mathrm{th}}(T)=\Phi_i^-(T)-\Phi_i^+(T),
\end{equation}
e una policy di Gibbs a temperatura $\Ta$ sceglie
\begin{equation}
  \pi_{\Ta}(x_i=\sigma)
  \propto\exp\bigl(-\Phi_i^\sigma/\Ta\bigr).
\end{equation}
Per il backtracking, il costo relativo dell'assegnazione corrente \`e
\begin{equation}
  \Delta\Phi_i
  =\Phi_i^{\mathrm{assigned}}-\min_{\sigma}\Phi_i^\sigma.
\end{equation}

\subsection{Survey virtuali e quantit\`a $I(k)$}

Un membro particolarmente semplice di questa famiglia si ottiene a vincolo duro ($\Tc=0$) senza liberare realmente la variabile.
Ricostruiamo le survey ``virtuali'' di ogni variabile fissata $k$ usando i messaggi correnti dei suoi vicini.
Se $k$ \`e fissata a vero, poniamo
\begin{equation}
  I(k)=1-s_F(k);
\end{equation}
se \`e fissata a falso,
\begin{equation}
  I(k)=1-s_T(k).
\end{equation}
Qui $s_T$ e $s_F$ sono le probabilit\`a survey di essere forzati a vero o a falso nella popolazione corrente dei cluster.
$I(k)$ \`e quindi la frazione di cluster attuali compatibili con il valore gi\`a assegnato a $k$.

Se i cluster della formula con $k$ fissata sono la frazione $I(k)$ dei cluster della formula in cui $k$ \`e libera, il guadagno di complessit\`a al rilascio vale
\begin{equation}
  \Delta\Sig_{\mathrm{release}}\simeq -\log I(k)
  \label{eq:delta-sigma-I}
\end{equation}
(logaritmo estensivo; nella convenzione densitaria il fattore $1/N$ compare esplicitamente).
Il backtracking dinamico-$I$ libera la variabile con il pi\`u piccolo $I(k)$ corrente.
Decimazione, direzione e frequenza $r$ restano quelle del BSP.

La relazione~\eqref{eq:delta-sigma-I} collega direttamente la teoria 1RSB al criterio operativo: liberare l'assegnazione compatibile con la frazione pi\`u piccola di cluster massimizza, a un passo, il recupero di $\Sig$.
Nella espansione a bassa temperatura $F_T=E_{\min}-T\log W_{\min}$, la differenza di free energy fratto $T$ riduce a un logaritmo di rapporti di probabilit\`a; $I(k)$ e $\Delta\Phi$ sono quindi due punti dello stesso funzionale, presi a temperature e a normalizzazioni diverse.

\subsection{Asse dei cluster non congelati}

Marino, Parisi e Ricci-Tersenghi~\cite{marino2016bsp} suggeriscono di concentrare la misura sui cluster unfrozen.
Una famiglia naturale \`e
\begin{equation}
  \mu_{m,\gamma}(C)
  \propto\exp\bigl[N(m\,s_C+\gamma\,q_C)\bigr],
\end{equation}
dove $q_C$ \`e la frazione di variabili non congelate del cluster.
Il parametro $m$ premia i cluster internamente ricchi; $\gamma$ premia i cluster in cui la variabile cavitarie non riceve warning.
Operativamente, nella stella a temperatura zero il settore unfrozen $A_0 B_0$ pesa $e^{\gamma}$ e i settori frozen restano invariati prima della rinormalizzazione.
I conflitti bilanciati a $\Tc>0$ ($p=q>0$) non ricevono $e^{\gamma}$.
Il messaggio resta $\eta=\pi^u/z^{(\gamma)}$.
Con $\gamma\neq 0$ la $\Sig$ stampata \`e la free entropy di questa misura biasata, non la complexity uniforme a $m=0$, $\gamma=0$.
L'opzione \`e \texttt{--rsb-gamma}. Non va confusa con lo scorer \texttt{gamma:<g>}, che ordina soltanto la decimazione.

\section{Metodi}
\label{sec:methods}

\subsection{Implementazione}

Le equazioni~\eqref{eq:F-thermo}--\eqref{eq:thermo-eta} sono implementate come somma polinomiale sui coefficienti dei prodotti $\prod_b(1-\eta_b+\eta_b z)$.
A $\Tc=0$ le correzioni di conflitto si annullano e i fattori riducono esattamente alle $\Pi$ del codice SP legacy.
La policy di Gibbs sulla direzione usa
\begin{equation}
  P(x_i=\mathrm{vero})
  =\bigl(1+e^{(\Phi_i^+-\Phi_i^-)/\Ta}\bigr)^{-1}
\end{equation}
e coincide con la regola dura per $\Ta\le 0$.
Il backtracking dinamico-$I$ ricostruisce le survey virtuali senza rilanciare una SP completa sulla formula con $k$ liberata.

\subsection{Validazione dello stimatore $I(k)$}

Su istanze intermedie del BSP, per $120$ scelte di variabili fissate, confrontiamo $-\log I(k)$ con la variazione reale di $\Sig$ ottenuta liberando $k$, riconvergendo SP e misurando $\Delta\Sig$.
Questa \`e la verifica meccanicistica centrale: lo stimatore locale deve predire il guadagno globale.

\subsection{Protocolli di risoluzione}

Confrontiamo tre criteri di backtracking a parit\`a di $r$ e di regola di decimazione:
\begin{itemize}
  \item certainty (BSP di riferimento);
  \item polarizzazione;
  \item dinamico-$I$.
\end{itemize}
I protocolli principali sono $K=3$, $N=10^4$, $\alpha=4{,}2$ (cinque seed) e $K=4$, $N=10^3$, $\alpha=9{,}5$ (dieci seed).
Oltre al tasso di successo SAT misuriamo la traiettoria $\Sig(t)$, la roughness a parit\`a di variabili fissate, il massimo drop di $\Sig$ a profondit\`a fissata e il tempo di calcolo.
Le soluzioni trovate sono verificate clausola per clausola sulla formula originale.

\section{Risultati}
\label{sec:results}

\subsection{Accuratezza di $I(k)$ come predittore di $\Delta\Sig$}

Nei $120$ esperimenti di liberazione reale, la correlazione fra guadagno previsto $-\log I(k)$ e guadagno misurato \`e $r=0{,}99684$.
Il segno della variazione \`e corretto in $120$ casi su $120$.
L'errore medio assoluto \`e $0{,}00056$.
Lo stimatore locale ricostruito dalle survey virtuali misura quindi, entro l'errore numerico, il guadagno globale di complessit\`a dovuto al rilascio.
Questo risultato supporta l'ipotesi che la risposta della free entropy al vincolo $x_k=$assegnato sia la quantit\`a rilevante per il backtracking, e non il bias congelato al momento del fissaggio.

\subsection{Dinamica della complessit\`a e tasso di successo}

Per $K=3$, $N=10^4$, $\alpha=4{,}2$, dinamico-$I$ risolve $5/5$ istanze, certainty $4/5$, polarizzazione $5/5$.
Rispetto a certainty, dinamico-$I$ riduce la discesa media di $\Sig$ del $26{,}0\%$, la roughness a parit\`a di variabili fissate del $17{,}7\%$ e il massimo drop a profondit\`a fissata del $17{,}9\%$, con un aumento di tempo del $5{,}8\%$.

Per $K=4$, $N=10^3$, $\alpha=9{,}5$, dinamico-$I$ risolve $6/10$ istanze, certainty $5/10$, polarizzazione $6/10$.
Rispetto a certainty, $|\Delta\Sig|$ medio scende dell'$8{,}6\%$, la discesa media del $23{,}2\%$, il massimo drop a profondit\`a fissata del $39{,}8\%$; la roughness migliora del $4{,}8\%$, mentre il tempo cresce di un fattore $1{,}94$.
Le undici soluzioni trovate da dinamico-$I$ sui due protocolli sono state verificate esaustivamente.

\subsection{Limiti osservati}

Una discesa pi\`u lenta di $\Sig$ non implica da sola lo spostamento della soglia algoritmica $\alpha_a$.
Un lookahead che massimizzava direttamente $\Sig$ dopo la mossa \`e risultato costosissimo e quasi inutile.
Alcuni criteri riducevano il massimo drop ma diminuivano il successo.
A $K=4$, $\alpha=9{,}7$, dinamico-$I$ produce ancora la curva di complessit\`a pi\`u regolare, ma tutti i metodi falliscono ($0/5$).
Conservare cluster disponibili \`e quindi una condizione favorevole, non sufficiente: contano anche la stabilit\`a del punto fisso SP e la scelta di assegnazioni che non conducano a futuri punti spinodali.

\section{Discussione}
\label{sec:discussion}

Il contributo pi\`u solido \`e meccanicistico: $I(k)$ collega una quantit\`a 1RSB locale alla variazione globale di $\Sig$, con accuratezza quasi esatta.
Questo risolve, nella pratica, il problema dei bias stantii nel backtracking del BSP.
La lettura in termini di free energy colloca $I(k)$ in una famiglia pi\`u ampia $(\Tc,\Ta,m,\gamma)$, di cui soltanto l'angolo duro della policy di rilascio \`e stato misurato in modo sistematico.

La deformazione ThermoSP a $\Tc>0$ coincide con SP a energia finita.
Il formalismo dei semianelli termodinamici fornisce linguaggio, forma variazionale e la domanda sulla conservazione della struttura di message passing sotto deformazioni di R\'enyi o Tsallis~\cite{marcolli2017entropy}.
Non fornisce da solo un nuovo algoritmo: senza un principio di selezione dei parametri, la famiglia rischia di diventare uno spazio di ricerca euristico.

Sull'asse $m>0$ siamo cauti.
I risultati con $I(k)$ stanno nel conteggio uniforme dei cluster ($m=0$).
Premiare i cluster a maggiore entropia interna allontana la misura dal punto sperimentalmente favorevole.
L'asse $\gamma$ verso cluster unfrozen resta, a nostro avviso, la direzione teorica pi\`u promettente e la meno sviluppata.

Il controllo soft globale basato su una ricorrenza di partition function sull'albero di ricerca~\cite{kappen2005path} richiede un troncamento.
Nei nostri test, i lookahead su $\Sig$ non hanno migliorato il successo a costo accettabile.
Preferiamo quindi politiche a un passo basate su risposte locali, come $I(k)$ e $\Delta\Phi$, rispetto a stime di valore a profondit\`a maggiore.

Per una rivendicazione sulla soglia algoritmica servono ancora una griglia di $\alpha$ vicino a $\alpha_s$, decine di seed per punto, stime di $\alpha_a$ con intervalli di confidenza, test a pi\`u $N$ e un insieme di formule tenute fuori dallo sviluppo.

\section{Conclusioni}
\label{sec:conclusions}

Abbiamo formulato la survey propagation termodinamica come deformazione di semianello dell'aggregazione cavitarie, con limite duro uguale a SP, e abbiamo collegato la policy di backtracking del BSP a risposte locali della free entropy.
Lo stimatore $I(k)$, ricostruito dalle survey virtuali, predice la variazione di complessit\`a al rilascio con correlazione prossima a uno e migliora la dinamica di $\Sig$ su istanze $3$-SAT e $4$-SAT di grandezza algoritmica rilevante.

Il passo successivo \`e derivare le equazioni SP con campo sorgente sulle variabili congelate (asse $\gamma$) e confrontare $\Delta\Phi$ a temperatura finita con $I(k)$ come predittore di $\Delta\Sig$ su formule non usate nello sviluppo.
Solo dopo quel confronto, e dopo una campagna statistica sulla soglia, sar\`a possibile decidere se la famiglia termodinamica sposta $\alpha_a$ oppure se il contributo pubblicabile resta la ricostruzione operativa delle survey virtuali.

\section*{Ringraziamenti}
Lavoro in corso.
Ringraziamo i colleghi e i collaboratori che hanno discusso le idee qui presentate.

\bibliography{thermosp}
\bibliographystyle{iclr2027_conference}

\end{document}
